\subsection{Vzdálenost dvou přímek}
\label{subsec:Distance_between_lines_space}

U dvou přímek v prostoru můžeme kromě jejich vzájemné polohy
a odchylky zkoumat také jejich vzdálenost.

\begin{definition}
    \label{def:Distance_between_lines_space}

    \textbf{Vzdáleností dvou přímek} $p$ a $q$ nazýváme nejmenší
    vzdálenost mezi libovolným bodem přímky $p$ a libovolným bodem
    přímky $q$.

    Platí tedy
    \[
        d(p,q)
        =
        \min_{\substack{X\in p\\Y\in q}} |XY|.
    \]
\end{definition}

Pokud mají přímky alespoň jeden společný bod, je jejich vzdálenost
rovna nule. Pro různoběžné přímky tedy platí
\[
    d(p,q)=0.
\]

Stejně tak je vzdálenost nulová u totožných přímek.

Nenulovou vzdálenost proto budeme určovat především ve dvou případech:
u přímek \textbf{rovnoběžných} a u přímek \textbf{mimoběžných}.


% ============================================================
%             R O V N O B Ě Ž N É   P Ř Í M K Y
% ============================================================

\subsubsection{Vzdálenost rovnoběžných přímek}
\label{subsubsec:Distance_parallel_lines_space}

Uvažujme dvě různé rovnoběžné přímky. Jejich vzdálenost odpovídá
délce úsečky, která spojuje obě přímky a je k nim kolmá.

Takových společných kolmic můžeme mezi dvěma rovnoběžnými přímkami
vést nekonečně mnoho. Všechny úsečky vymezené jejich průsečíky
s oběma přímkami však mají stejnou délku.

Nejprve si ukážeme výpočet přímo podle této geometrické představy.

\begin{example}
    \label{exa:Distance_parallel_lines_geometrically}

    Určete vzdálenost rovnoběžných přímek
    \[
        \begin{aligned}
            p: x&=t,\\
            y&=t,\\
            z&=0; t\in\mathbb{R}
        \end{aligned}
    \]
    a
    \[
        \begin{aligned}
            q: x&=s,\\
            y&=2+s,\\
            z&=1; s\in\mathbb{R}.
        \end{aligned}
    \]

    Směrové vektory obou přímek jsou
    \[
        \vec{u}=(1;1;0),
    \]
    takže jsou přímky rovnoběžné.

    Chceme najít přímku $r$, která protíná přímky $p$ a $q$
    a je k nim kolmá.

    Na přímce $p$ si zvolíme například bod
    \[
        A[0;0;0].
    \]

    Druhý průsečík společné kolmice s přímkou $q$ označíme $B$.
    Protože $B\in q$, můžeme jeho souřadnice zapsat pomocí parametru
    $s$:
    \[
        B[s;2+s;1].
    \]

    Vektor spojující body $A$ a $B$ je
    \[
        \overrightarrow{AB}
        =
        (s;2+s;1).
    \]

    Přímka $AB$ má být kolmá k přímce $p$, a proto musí být
    vektor $\overrightarrow{AB}$ kolmý ke směrovému vektoru
    $\vec{u}$.

    Musí tedy platit
    \[
        \overrightarrow{AB}\cdot\vec{u}=0.
    \]

    Po dosazení dostáváme
    \begin{align*}
        (s;2+s;1)\cdot(1;1;0)&=0,\\
        s+2+s&=0,\\
        2s+2&=0,\\
        s&=-1.
    \end{align*}

    Hledaný bod na přímce $q$ je tedy
    \[
        B[-1;1;1].
    \]

    Směrovým vektorem společné kolmice může být
    \[
        \overrightarrow{AB}=(-1;1;1).
    \]

    Přímku $r$ tedy můžeme zapsat parametricky ve tvaru
    \[
        \begin{aligned}
            r: x&=-\lambda,\\
            y&=\lambda,\\
            z&=\lambda; \lambda\in\mathbb{R}.
        \end{aligned}
    \]

    Pro $\lambda=0$ dostáváme bod
    \[
        A[0;0;0]\in p,
    \]
    zatímco pro $\lambda=1$ dostáváme bod
    \[
        B[-1;1;1]\in q.
    \]

    Navíc platí
    \[
        (-1;1;1)\cdot(1;1;0)=0,
    \]
    takže přímka $r$ je skutečně kolmá k oběma rovnoběžkám.

    Vzdálenost přímek $p$ a $q$ je proto již pouze vzdáleností
    bodů $A$ a $B$:
    \begin{align*}
        d(p,q)
        &=
        |AB|\\
        &=
        \sqrt{(-1-0)^2+(1-0)^2+(1-0)^2}\\
        &=
        \sqrt{3}.
    \end{align*}

    Platí tedy
    \[
        \boxed{d(p,q)=\sqrt{3}}.
    \]
\end{example}

Předchozí postup přesně odpovídá geometrické představě vzdálenosti
rovnoběžných přímek. Nejprve jsme našli jejich společnou kolmici,
poté její průsečíky s oběma přímkami a nakonec jsme již pouze
vypočítali vzdálenost těchto dvou bodů.

Pokud nás však zajímá pouze vzdálenost přímek a nepotřebujeme znát
společnou kolmici ani její průsečíky, můžeme celý výpočet výrazně
zjednodušit.

Uvažujme dvě různé rovnoběžné přímky
\[
    p:\ X=A+t\vec{u}
    \quad\text{a}\quad
    q:\ X=B+s\vec{u},
\]
kde $A\in p$, $B\in q$ a $\vec{u}$ je jejich společný směrový vektor.

Protože jsou přímky rovnoběžné, mají všechny body jedné přímky
od druhé přímky stejnou vzdálenost. Stačí tedy zvolit libovolný
bod jedné z přímek a určit jeho vzdálenost od druhé přímky.

Například zvolíme bod $B\in q$. Potom
\[
    d(p,q)=d(B,p).
\]

Pro vzdálenost bodu od přímky jsme již odvodili vztah
\[
    d(B,p)
    =
    \frac{
        |\overrightarrow{AB}\times\vec{u}|
    }{
        |\vec{u}|
    }.
\]

Proto pro dvě rovnoběžné přímky platí

\begin{equation}
    \boxed{
        d(p,q)
        =
        \frac{
            |\overrightarrow{AB}\times\vec{u}|
        }{
            |\vec{u}|
        }
    }
    \label{eq:Distance_parallel_lines_space}
\end{equation}

kde $A$ je libovolný bod přímky $p$, $B$ libovolný bod přímky $q$
a $\vec{u}$ směrový vektor obou přímek.

\begin{example}
    \label{exa:Distance_parallel_lines_formula_space}

    Vraťme se k přímkám z předchozího příkladu:
    \[
        \begin{aligned}
            p: x&=t,\\
            y&=t,\\
            z&=0
        \end{aligned}
        \qquad\text{a}\qquad
        \begin{aligned}
            q: x&=s,\\
            y&=2+s,\\
            z&=1.
        \end{aligned}
    \]

    Místo hledání společné kolmice tentokrát zvolíme libovolné body
    obou přímek, například
    \[
        A[0;0;0]\in p,
        \qquad
        B[0;2;1]\in q.
    \]

    Potom
    \[
        \overrightarrow{AB}=(0;2;1)
    \]
    a společný směrový vektor přímek je
    \[
        \vec{u}=(1;1;0).
    \]

    Vektorový součin je
    \[
        \overrightarrow{AB}\times\vec{u}
        =
        (0;2;1)\times(1;1;0)
        =
        (-1;1;-2).
    \]

    Proto
    \[
        |\overrightarrow{AB}\times\vec{u}|
        =
        \sqrt{1+1+4}
        =
        \sqrt6
    \]
    a
    \[
        |\vec{u}|=\sqrt2.
    \]

    Dosadíme:
    \[
        d(p,q)
        =
        \frac{\sqrt6}{\sqrt2}
        =
        \boxed{\sqrt3}.
    \]

    Dostáváme stejnou vzdálenost jako předchozím geometrickým
    postupem, tentokrát však bez hledání společné kolmice.
\end{example}

\begin{remark}
    Při výpočtu můžeme zaměnit přímky $p$ a $q$, protože
    \[
        d(p,q)=d(q,p).
    \]

    Stejně tak můžeme použít libovolný nenulový násobek směrového
    vektoru. Výsledná vzdálenost se tím nezmění.
\end{remark}

\begin{example}
    \label{exa:Find_parallel_lines_given_distance_space}

    Je dána přímka
    \[
        \begin{aligned}
            p: x&=1+t,\\
            y&=-1+t,\\
            z&=2; t\in\mathbb{R}.
        \end{aligned}
    \]

    Určete všechny přímky $q$, které jsou rovnoběžné s přímkou $p$,
    protínají osu $z$ a jejich vzdálenost od přímky $p$ je rovna $3$.

    Směrový vektor přímky $p$ je
    \[
        \vec{u}=(1;1;0).
    \]

    Protože má být přímka $q$ rovnoběžná s přímkou $p$, můžeme pro ni
    použít stejný směrový vektor.

    Přímka $q$ má zároveň protínat osu $z$. Její průsečík s osou $z$
    proto můžeme označit
    \[
        B[0;0;m],
    \]
    kde $m\in\mathbb{R}$.

    Přímka $q$ má tedy parametrické vyjádření
    \[
        \begin{aligned}
            q: x&=s,\\
            y&=s,\\
            z&=m; s\in\mathbb{R}.
        \end{aligned}
    \]

    Z přímky $p$ zvolíme například bod
    \[
        A[1;-1;2].
    \]

    Potom
    \[
        \overrightarrow{AB}
        =
        B-A
        =
        (-1;1;m-2).
    \]

    Vzdálenost rovnoběžných přímek určíme ze vztahu
    \[
        d(p,q)
        =
        \frac{
            |\overrightarrow{AB}\times\vec{u}|
        }{
            |\vec{u}|
        }.
    \]

    Nejprve vypočítáme vektorový součin:
    \begin{align*}
        \overrightarrow{AB}\times\vec{u}
        &=
        (-1;1;m-2)\times(1;1;0)\\
        &=
        (2-m;m-2;-2).
    \end{align*}

    Jeho velikost je
    \begin{align*}
        |\overrightarrow{AB}\times\vec{u}|
        &=
        \sqrt{(2-m)^2+(m-2)^2+(-2)^2}\\
        &=
        \sqrt{2(m-2)^2+4}.
    \end{align*}

    Dále
    \[
        |\vec{u}|=\sqrt2.
    \]

    Proto
    \[
        d(p,q)
        =
        \frac{\sqrt{2(m-2)^2+4}}{\sqrt2}
        =
        \sqrt{(m-2)^2+2}.
    \]

    Požadujeme
    \[
        d(p,q)=3.
    \]

    Dostáváme
    \begin{align*}
        \sqrt{(m-2)^2+2}&=3,\\
        (m-2)^2+2&=9,\\
        (m-2)^2&=7,\\
        m-2&=\pm\sqrt7,
    \end{align*}
    a tedy
    \[
        m=2\pm\sqrt7.
    \]

    Existují proto dvě hledané přímky:
    \[
        \begin{aligned}
            q_1: x&=s,\\
                 y&=s,\\
                 z&=2+\sqrt7; s\in\mathbb{R},
        \end{aligned}
        \qquad
        \begin{aligned}
            q_2: x&=s,\\
                 y&=s,\\
                 z&=2-\sqrt7; s\in\mathbb{R}.
        \end{aligned}
    \]

    Pro obě platí
    \[
        d(p,q_1)=d(p,q_2)=3.
    \]
\end{example}


% ============================================================
%             M I M O B Ě Ž N É   P Ř Í M K Y
% ============================================================

\subsubsection{Vzdálenost mimoběžných přímek}
\label{subsubsec:Distance_skew_lines_space}

U rovnoběžných přímek jsme mohli zvolit bod jedné přímky a hledat
jeho kolmý průmět na přímku druhou. U mimoběžek je situace
složitější, protože jejich směry nejsou rovnoběžné.

Také zde však vzdálenost odpovídá délce nejkratší úsečky spojující
obě přímky. Tato úsečka musí být kolmá současně k oběma přímkám.

\begin{definition}
    \label{def:Common_perpendicular_skew_lines}

    Přímku, která protíná dvě mimoběžné přímky a je kolmá k oběma
    z nich, nazýváme jejich \textbf{společnou kolmicí}.

    Úsečka společné kolmice vymezená jejími průsečíky s oběma
    mimoběžkami má nejmenší možnou délku ze všech úseček spojujících
    body těchto přímek. Její délka je vzdáleností mimoběžných přímek.
\end{definition}

Nejprve opět určíme vzdálenost přímo podle této geometrické
představy.

\begin{example}
    \label{exa:Distance_skew_lines_geometrically}

    Určete vzdálenost mimoběžných přímek
    \[
        \begin{aligned}
            p: x&=t,\\
            y&=t,\\
            z&=0; t\in\mathbb{R}
        \end{aligned}
    \]
    a
    \[
        \begin{aligned}
            q: x&=s,\\
            y&=0,\\
            z&=-3+s; s\in\mathbb{R}.
        \end{aligned}
    \]

    Směrové vektory přímek jsou
    \[
        \vec{u}=(1;1;0),
        \qquad
        \vec{v}=(1;0;1).
    \]

    Vektory nejsou rovnoběžné a přímky nemají společný bod, takže
    jsou přímky mimoběžné.

    Hledáme jejich společnou kolmici. Její průsečík s přímkou $p$
    označíme $P$ a průsečík s přímkou $q$ označíme $Q$.

    Protože $P\in p$, můžeme psát
    \[
        P[t;t;0].
    \]

    Podobně pro $Q\in q$ platí
    \[
        Q[s;0;-3+s].
    \]

    Vektor spojující oba body je
    \begin{align*}
        \overrightarrow{PQ}
        &=
        Q-P\\
        &=
        (s-t;-t;-3+s).
    \end{align*}

    Společná kolmice musí být kolmá k přímce $p$. Proto
    \[
        \overrightarrow{PQ}\cdot\vec{u}=0.
    \]

    Dosadíme:
    \begin{align*}
        (s-t;-t;-3+s)\cdot(1;1;0)&=0,\\
        s-2t&=0.
    \end{align*}

    Zároveň musí být společná kolmice kolmá k přímce $q$, a proto
    \[
        \overrightarrow{PQ}\cdot\vec{v}=0.
    \]

    Dostáváme
    \begin{align*}
        (s-t;-t;-3+s)\cdot(1;0;1)&=0,\\
        s-t-3+s&=0,\\
        2s-t-3&=0.
    \end{align*}

    Hledané parametry tedy splňují soustavu
    \[
        \begin{cases}
            s-2t=0,\\
            2s-t-3=0.
        \end{cases}
    \]

    Z první rovnice
    \[
        s=2t.
    \]

    Dosadíme do druhé:
    \begin{align*}
        2(2t)-t-3&=0,\\
        3t&=3,\\
        t&=1.
    \end{align*}

    Odtud
    \[
        s=2.
    \]

    Průsečíky společné kolmice s danými mimoběžkami jsou tedy
    \[
        P[1;1;0]
    \]
    a
    \[
        Q[2;0;-1].
    \]

    Vektor
    \[
        \overrightarrow{PQ}
        =
        (1;-1;-1)
    \]
    je skutečně kolmý k oběma směrovým vektorům:
    \[
        (1;-1;-1)\cdot(1;1;0)=0
    \]
    a
    \[
        (1;-1;-1)\cdot(1;0;1)=0.
    \]

    Společná kolmice může být tedy zapsána například ve tvaru
    \[
        \begin{aligned}
            r: x&=1+\lambda,\\
            y&=1-\lambda,\\
            z&=-\lambda; \lambda\in\mathbb{R}.
        \end{aligned}
    \]

    Pro $\lambda=0$ získáme bod $P$ a pro $\lambda=1$ bod $Q$.

    Nyní již stačí určit vzdálenost těchto dvou bodů:
    \begin{align*}
        d(p,q)
        &=
        |PQ|\\
        &=
        \sqrt{(2-1)^2+(0-1)^2+(-1-0)^2}\\
        &=
        \sqrt3.
    \end{align*}

    Platí tedy
    \[
        \boxed{d(p,q)=\sqrt3}.
    \]
\end{example}

Předchozí postup přesně odpovídá geometrickému významu vzdálenosti
mimoběžných přímek. Museli jsme však najít dva neznámé body,
sestavit soustavu rovnic, určit společnou kolmici a teprve potom
vypočítat vzdálenost jejích průsečíků s oběma přímkami.

Pokud potřebujeme pouze samotnou vzdálenost, můžeme postup výrazně
zkrátit.

Uvažujme mimoběžné přímky
\[
    p:\ X=A+t\vec{u}
\]
a
\[
    q:\ X=B+s\vec{v}.
\]

Protože nejsou jejich směrové vektory rovnoběžné, platí
\[
    \vec{u}\times\vec{v}\neq\vec{0}.
\]

Vektor
\[
    \vec{u}\times\vec{v}
\]
je kolmý současně ke směrovým vektorům obou přímek. Jeho směr je
tedy rovnoběžný se směrem společné kolmice mimoběžek.

Zvolíme libovolné body
\[
    A\in p,
    \qquad
    B\in q.
\]

Vektory
\[
    \vec{u},\qquad
    \vec{v},\qquad
    \overrightarrow{AB}
\]
určují rovnoběžnostěn.

Jeho objem můžeme pomocí smíšeného součinu vyjádřit jako
\[
    V
    =
    \left|
        \overrightarrow{AB}
        \cdot
        (\vec{u}\times\vec{v})
    \right|.
\]

Obsah podstavy tohoto rovnoběžnostěnu určené vektory
$\vec{u}$ a $\vec{v}$ je
\[
    S
    =
    |\vec{u}\times\vec{v}|.
\]

Výška rovnoběžnostěnu je právě vzdáleností mimoběžných přímek
$d(p,q)$. Proto platí také
\[
    V
    =
    |\vec{u}\times\vec{v}|\cdot d(p,q).
\]

Porovnáním obou vyjádření objemu dostáváme
\[
    |\vec{u}\times\vec{v}|\cdot d(p,q)
    =
    \left|
        \overrightarrow{AB}
        \cdot
        (\vec{u}\times\vec{v})
    \right|.
\]

Odtud

\begin{equation}
    \boxed{
        d(p,q)
        =
        \frac{
            \left|
                \overrightarrow{AB}
                \cdot
                (\vec{u}\times\vec{v})
            \right|
        }{
            |\vec{u}\times\vec{v}|
        }
    }
    \label{eq:Distance_skew_lines_space}
\end{equation}

kde $A\in p$ a $B\in q$ jsou libovolné body daných mimoběžek
a $\vec{u}$, $\vec{v}$ jsou jejich směrové vektory.

\begin{example}
    \label{exa:Distance_skew_lines_formula_space}

    Vraťme se k mimoběžkám z předchozího příkladu:
    \[
        \begin{aligned}
            p: x&=t,\\
            y&=t,\\
            z&=0
        \end{aligned}
        \qquad\text{a}\qquad
        \begin{aligned}
            q: x&=s,\\
            y&=0,\\
            z&=-3+s.
        \end{aligned}
    \]

    Zvolíme libovolné body
    \[
        A[0;0;0]\in p,
        \qquad
        B[0;0;-3]\in q.
    \]

    Potom
    \[
        \overrightarrow{AB}=(0;0;-3).
    \]

    Směrové vektory přímek jsou
    \[
        \vec{u}=(1;1;0),
        \qquad
        \vec{v}=(1;0;1).
    \]

    Jejich vektorový součin je
    \[
        \vec{u}\times\vec{v}
        =
        (1;-1;-1).
    \]

    Proto
    \[
        |\vec{u}\times\vec{v}|=\sqrt3.
    \]

    Dále
    \begin{align*}
        \overrightarrow{AB}\cdot(\vec{u}\times\vec{v})
        &=
        (0;0;-3)\cdot(1;-1;-1)\\
        &=
        3.
    \end{align*}

    Dosadíme do vztahu pro vzdálenost mimoběžek:
    \[
        d(p,q)
        =
        \frac{|3|}{\sqrt3}
        =
        \boxed{\sqrt3}.
    \]

    Dostáváme stejný výsledek jako při hledání společné kolmice,
    tentokrát však bez nutnosti určovat její průsečíky s oběma
    přímkami.
\end{example}

\begin{remark}
    Vztah
    \[
        d(p,q)
        =
        \frac{
            \left|
                \overrightarrow{AB}
                \cdot
                (\vec{u}\times\vec{v})
            \right|
        }{
            |\vec{u}\times\vec{v}|
        }
    \]
    nemůžeme použít pro rovnoběžné přímky.

    Pokud by byly jejich směrové vektory rovnoběžné, platilo by
    \[
        \vec{u}\times\vec{v}=\vec{0},
    \]
    a jmenovatel vztahu by byl nulový.
\end{remark}

Vzdálenost dvou přímek v prostoru tedy závisí na jejich vzájemné
poloze:

\[
    \begin{array}{c|c}
        \text{vzájemná poloha přímek} & \text{vzdálenost}\\
        \hline
        \text{totožné} & 0\\
        \text{různoběžné} & 0\\
        \text{rovnoběžné}
        &
        \displaystyle
        \frac{|\overrightarrow{AB}\times\vec{u}|}{|\vec{u}|}\\[3mm]
        \text{mimoběžné}
        &
        \displaystyle
        \frac{
            \left|
                \overrightarrow{AB}\cdot
                (\vec{u}\times\vec{v})
            \right|
        }{
            |\vec{u}\times\vec{v}|
        }
    \end{array}
\]